%% ----------------------------------------------------------------
%   Template created by Dr. Malay Bhattacharyya, Department of IT
%                           IIEST, Shibpur
%                           Copyright 2015
%% ----------------------------------------------------------------
\documentclass[12pt]{article}
% Packages to be added here
\usepackage{geometry}
\geometry{verbose,a4paper,tmargin=30mm,bmargin=30mm,lmargin=30mm,rmargin=20mm}
\usepackage{graphics}
\usepackage{setspace}
\usepackage{url}
\usepackage{epsfig}
\usepackage{amsmath}
\usepackage{amsthm}
\usepackage{amssymb}
\usepackage{multirow}
\usepackage{multicol}
\usepackage{algorithm}
\usepackage{algorithmic}
\usepackage{hyperref}

% For definitions, theorems, lemmas, proofs, corollaries and observations
\newtheorem{DEF}{Definition}
\newtheorem{THM}{Theorem}[section]
\newtheorem{LEM}{Lemma}[section]
\newtheorem{COR}{Corollary}[section]
\newtheorem{PRF}{Proof}[]
\newtheorem{OBS}{Observation}[section]

\setcounter{secnumdepth}{5}
\setcounter{tocdepth}{5}

\onehalfspacing

\makeatletter

%%%%%%%%%%%%%%%%%%%%%%%%%%%%%% LyX specific LaTeX commands.
\newcommand{\LyX}{L\kern-.1667em\lower.25em\hbox{Y}\kern-.125emX\spacefactor1000}

\makeatother

\begin{document}




% Title page
\begin{titlepage}
\thispagestyle{empty}
{\centering
{\large A Project Report on}\\
\vspace{0.5cm}
\large
{\Huge ALGORITHMIC ART USING RANDOM ISOTHETIC POLYGON}\\
\vspace{0.5cm}
Submitted to the Department of Information Technology\\
\vspace{1cm}
{\Large\bf For the partial fulfilment of the degree of B.E. in Information Technology}\\
\vspace{1cm}
by \\
\vspace{.5cm}
\hspace{.05cm} {\large \bf{Binayak Ghosh}}\\
\hspace{.05cm} {\large \bf{Ronit Halder}}\\
\hspace{.05cm} {\large \bf{Firoz Rahaman}}\\
\vspace {0.3cm}
\hspace{.05cm} Under the supervision of\\
\hspace{.05cm} {\large \bf{Arindam Biswas}}\\
\vspace {0.8cm}
Roll number: Hx02, Hy34, Hx01\\
Registration number: 110812005, 110812025, 110812047 of 2013-14\\
B.E. 3rd year\\
\vspace*{0.7cm}
\begin{figure}[h]
{\centering {\includegraphics[scale=1.5]{Logo_IIEST.eps}}\par}
\end{figure}
Department of Information Technology\\
INDIAN INSTITUTE OF ENGINEERING SCIENCE AND TECHNOLOGY, SHIBPUR\\
\begin{center}
\textit{May, 2015}
\end{center}

{\centering
\hspace{6.5cm} \date{\today}}
}
\pagebreak
\end{titlepage}




\newpage




% Certificate page
\thispagestyle{empty}
\begin{center}
  \large \textbf{Department of Information Technology\\
  Indian Institute of Engineering Science and Technology, Shibpur}
\begin{figure}
  \centering
  \includegraphics[scale=1]{Logo_IIEST.eps}\\
\end{figure}
\\
 \vspace{1cm}
 \huge \textbf{CERTIFICATE}
 \vspace{1cm}
\end{center}
This is to certify that the work presented in this report entitled ``ALGORITHMIC ART USING RANDOM ISOTHETIC POLYGONS'', submitted by Binayak Ghosh, Ronit Halder and Firoz Rahaman having the examination roll numbers 111208005, 111208025 and 111208047 respectively has been carried out under my supervision for the partial fulfilment of the degree of Bachelor of Technology in Information Technology during the session 2014-15 in the Department of Information Technology, Indian Institute of Engineering Science and Technology, Shibpur.\\
\vspace{1.5cm}

\begin{minipage}[t]{0.48\linewidth}
\begin{flushleft}
---------------------------------------\\
Arindam Biswas\\
Associate Professor\\
Department of Information Technology\\
Indian Institute of Engineering Science and Technology, Shibpur
\end{flushleft}
\end{minipage}\hfill
\begin{minipage}[t]{0.48\linewidth}
\begin{flushright}
---------------------------------------\\
Arindam Biswas\\
Head of the Department\\
Department of Information Technology\\
Indian Institute of Engineering Science and Technology, Shibpur
\end{flushright}
\end{minipage}
\vspace{1.5cm}

\begin{minipage}[t]{0.48\linewidth}
-----------------------------------------\\
Dean (Academic)
\end{minipage}\hfill
\begin{minipage}[t]{0.48\linewidth}
\begin{center}
Date: 08/05/2015
\end{center}
\end{minipage}




\newpage




% Acknowledgment page
\thispagestyle{empty}
\begin{center}
  \Huge \textbf{Acknowledgements}
\end{center}
\vspace{2cm}
Being a novice in the area of research, this project would not have been possible without constant
supervision of Prof. Arindam Biswas. The combinatorial algorithm for construction of isothetic cover of digital objects, developed by Prof. Arindam Biswas and his colleagues, formed the basis of our project. Also, the works of Mike Bostock in the field of algorithmic visualisation have served as great source for inspiration.\\
\vspace{2.5cm}

\begin{minipage}[t]{0.48\linewidth}
\begin{flushleft}
Date: 08/05/2015
\end{flushleft}
\end{minipage}\hfill
\begin{minipage}[t]{0.48\linewidth}
\begin{flushright}
---------------------------------------\\
Binayak Ghosh,\\
Ronit Halder,\\
Firoz Rahaman\\
Department of Information Technology\\
Indian Institute of Engineering Science and Technology, Shibpur
\end{flushright}
\end{minipage}




\newpage




% Abstract page
\begin{abstract}
This report summarizes the work done as part of a semester project under the supervision of Prof. Arindam Biswas. The project explores the fascinating field of generative art. Using random isothetic polygons interesting visually stimulating patterns have been generated from existing images. These patterns are shown in the report and the processes used to generate them have been discussed.

\end{abstract}




% Other front pages
\pagenumbering{roman}
\newpage
\tableofcontents
\newpage
\newenvironment{mydef}[1]{\begin{definition} #1 \mbox{\\}
\rm}{\end{definition}}
\pagenumbering{arabic}




\section{INTRODUCTION}
Algorithmic art or generative art refers to art (mainly visual art) that in whole or in part has been created with the use of an autonomous system. In such an art form the generative system is said to be the creation of the artist which in turn creates a set of images which are unique and visually appealing.  
So generation of patterns a system needs to be created which on being provided some form of input will produce a set of images.

\vspace{0.3cm} 
Isothetic polygons can be used to generate isothetic covers of binary images which can then be divided into unit gird blocks. These blocks can then be filled with different colors based on some algorithms to produce visually appealing patterned images. Thus using isothetic polygons a generation system can be created which can produce some random patterns based on binary images of regular day to day objects/shapes. 
\vspace{0.3cm}





\newpage





\section{PRELIMINARIES AND DEFINITIONS}
\begin{DEF}
  \textnormal{Digital Plane}: The digital plane, $\textbf{Z}^2$, is the set of all points having integer coordinates in the real plane $\textbf{R}^2$.A point in the digital plane is called a digital point,or called a pixel in case of a digital image. Henceforth, the term "point" and "digital point" will be interchangeably used.
\end{DEF}


\begin{DEF}
\textnormal{Digital grid}: A grid \textbf{G} := (\textbf{H},\textbf{V}) consists of a set \textbf{H} of horizontal grid lines and a set \textbf{V} of vertical grid lines, where, \textbf{H} = \{...,${\textit{l}}_H$(\textit{j-2g}),${\textit{l}}_H$(\textit{j-g}), ${\textit{l}}_H$(\textit{j}), ${\textit{l}}_H$(\textit{j+g}), ${\textit{l}}_H$(\textit{j+2g}),...\}{$\subset $$\textbf{Z}^2$} and \textbf{V}=\{...,${\textit{l}}_V$(\textit{j-2g}),${\textit{l}}_V$(\textit{j-g}), ${\textit{l}}_V$(\textit{j}), ${\textit{l}}_V$(\textit{j+g}), ${\textit{l}}_V$(\textit{j+2g}),...\}{$\subset $$\textbf{Z}^2$}, for a grid size,\textit{g} $\in$ {$\textbf{Z}^+$}. Here, ${\textit{l}}_H$ := \{(\textit{ $\i^{'}$},\textit{j}):\textit{$\i^{'}$}$\in$ \textbf{Z} denotes the horizontal grid line and ${\textit{l}}_V$ := \{(\textit{i},\textit{ $\j^{'}$}):\textit{$\j^{'}$}$\in$ \textbf{Z} denotes the vertical grid line intersecting at point (\textit{i},\textit{j})$\in$ $\textbf{Z}^2$,called the grid point,where i and j multiples of g.
\end{DEF}

\begin{DEF}
\textnormal{Grid segment}: A grid segment is the set of points on a grid line between and inclusive of two consecutive grid point.Thus,the horizontal grid segment belonging to ${\textit{l}}_H$(\textit{j}) and lying between two consecutive vertical grid lines, ${\textit{l}}_V$(\textit{i}) and ${\textit{l}}_V$(\textit{i}+\textit{j}),is  \textit{$s_j^{i}$} =\{(\textit{ $\i^{'}$},\textit{j})$\in$ ${\textit{l}}_H$(\textit{j}):\textit{i}$\leq$\textit{$i^{'}$}$\leq$\textit{i}+\textit{g}\}.Similarly,the vertical the vertical grid segment belonging to ${\textit{l}}_V$(\textit{i}) between ${\textit{l}}_H$(\textit{j}) and ${\textit{l}}_V$(\textit{i}+\textit{j}), is \textit{$s_i^{j}$} =\{(\textit{i},\textit{ $\j^{'}$})$\in$ ${\textit{l}}_H$(\textit{j}):\textit{j}$\leq$\textit{$j^{'}$}$\leq$\textit{j}+\textit{g}\}.
\end{DEF}


\begin{DEF}
\textnormal{Unit grid block (UGB)}: The horizontal digital line ${\textit{l}}_H$(\textit{j}) divides $\textbf{Z}^2$ into half-planes ,which are given by  \textit{$h_H^{+}$}(\textit{j}):=\{(\textit{ $\i^{'}$},\textit{ $\j^{'}$})$\in$ $\textbf{Z}^2$:\textit{$\j^{'}$}$\geq$\textit{j}\} and  \textit{$h_H^{-}$}(\textit{j}):=\{(\textit{ $\i^{'}$},\textit{ $\j^{'}$})$\in$ $\textbf{Z}^2$:\textit{$\j^{'}$}$\leq$\textit{j}\}. Similarly, ${\textit{l}}_V$(\textit{i}) divides \textbf{$Z^{2}$} into two half-planes, which are \textit{$h_V^{+}$}(\textit{i}):=\{(\textit{ $\i^{'}$},\textit{ $\j^{'}$})$\in$ $\textbf{Z}^2$:\textit{$\i^{'}$}$\geq$\textit{i}\} and  \textit{$h_V^{-}$}(\textit{j}):=\{(\textit{ $\i^{'}$},\textit{ $\j^{'}$})$\in$ $\textbf{Z}^2$:\textit{$\i^{'}$}$\leq$\textit{i}\}. Then the set, given by (\textit{$h_V^{+}$}(\textit{i})$\cap$\textit{$h_V^{-}$}(\textit{i}+\textit{g})$\cap$\textit{$h_H^{+}$}(\textit{j})$\cap$\textit{$h_H^{-}$}(\textit{j}+\textit{g})),is defined as the unit grid block,UGB(\textit{i},\textit{j}).The interior of \textit{$U_{1}$} :=UGB(\textit{i},\textit{j}) is given by \textit{$U_{1}$} \ (\textit{$s_{i}^{j}$}$\cup$\textit{$s_{i+g}^{j}$}$\cup$ \textit{$s_{i}^{j}$}$\cup$\textit{$s_{i+g}^{j}$}). Hence, for a given grid point, qði; jÞ, there are four neighbouring UGBs, namely, \textit{$U_{1}$} :=UGB(\textit{i},\textit{j}),\textit{$U_{2}$} :=UGB(\textit{i-g},\textit{j}),\textit{$U_{3}$} :=UGB(\textit{i-g},\textit{j-g}) and \textit{$U_{4}$} :=UGB(\textit{i-g},\textit{j}).
\end{DEF}
 

\begin{DEF}
\textnormal{Isothetic polygon}: An isothetic polygon \textit{P} is a simple polygon of finite size in \textbf{$Z^{2}$} whose alternative sides are subsets of the members of \textbf{H} and \textbf{V}.The polygon \textit{P},hence given by a finite set of UGBs,is represented by the sequence of its vertices,which are grid points.The border \textit{$B_{P}$} of \textit{P} is the set of points belonging to its sides.The interior of \textit{P} is the set of points in the union of its constituting UGBs excluding the border of P.
\end{DEF}


\begin{DEF}
\textnormal{Outer isothetic cover}: The outer (isothetic) cover (OIC), denoted by $\bar{P}$(\textit{S}), is a set of outer polygons and (outer) hole polygons, such that the region, given by the union of the outer polygons minus the union of the interiors of the hole polygons, contains a UGB if and only if it has object occupancy (i.e., has a nonempty intersection with S).
\end{DEF}


\newpage




\section{PROBLEM DEFINITION}
To generate outer isothetic cover from a given input image. Divide the isothetic cover into Unit Grid Blocks and then fill each UGB with a color based on some algorithms thereby generate visually appealing patterns.


\section{PROCESSES USED}
\subsection{Genarate Isothetic Cover}
We used the combinatorial algorithm developed by Arindam Biswas, Partha Bhowmick and Bhargab B. Bhattacharya as stated in their paper ``Construction of isothetic covers of a digital object: A combinatorial approach"\textsuperscript{\cite{lamport94}} to construct an outer isothetic cover from a given binary image. We then divided the isothetic cover into Unit Grid Blocks (UGBs) and animated the steps to create a nice visualisation. The following figure show different phases of the animation: 


\subsection{Fill UGBs with Random Color}

\begin{algorithm}[htp]
\caption{blurr-effect}
\textbf{Input:} I(The image as 2D Matrix),gsize(Grid size).\par 
\textbf{Output:} Modified Image I with blur effect.\par
\textbf{Steps of the algorithm:}
\begin{algorithmic}[1]
\STATE ncols $\longleftarrow$ width of input image I
\STATE nrows $\longleftarrow$ hight of input image I
\STATE start $\longleftarrow$ 4/5 part of gsize
\FOR{i $\longleftarrow$ 1 to nrows}
	\STATE P $\longleftarrow$ ith row of image as a list of pixels
	\IF{j th element of P is not a white pixel}
		\STATE count $\longleftarrow$ 0
		\FOR{i $\longleftarrow$ (j+start) to (j+gsize-1)}		
			\STATE count $\longleftarrow$ count+1
			\STATE $P_k$$\longleftarrow$$P_k$$\times$((2$\times$(gsize-start)-count)/2$\times$(gc-start)+((cnt/2)$\times$(gc-start))$\times$ $P_{(k+gsize)}$
		\ENDFOR
		\FOR{i $\longleftarrow$ (j+1) to (j+2$\times$gsize)}		
			\STATE count $\longleftarrow$ count+1
			\STATE $P_{k}$$\longleftarrow$$P_{(k-gsize)}$$\times$((2$\times$(gsize-start)-count)/2$\times$(gc-start)+((cnt/2)$\times$(gc-start))$\times$ $P_{k}$
		\ENDFOR
   \ENDIF
\ENDFOR
\FOR{i $\longleftarrow$ 1 to ncols}
	\STATE P $\longleftarrow$ $i^{th}$ column of image as a list of pixels
	\IF{$i^{th}$ element of P is not a white pixel}
		\STATE count $\longleftarrow$ 0
		\FOR{j $\longleftarrow$ (i+start+1) to (i+gsize)}		
			\STATE count $\longleftarrow$ count+1
			\STATE $P_k$$\longleftarrow$$P_k$$\times$((2$\times$(gsize-start)-count)/2$\times$(gc-start)+((cnt/2)$\times$(gc-start))$\times$ $P_{(k+gsize)}$
		\ENDFOR
		\FOR{j $\longleftarrow$ (i+1) to (i+2$\times$gsize)}		
			\STATE count $\longleftarrow$ count+1
			\STATE $P_{k}$$\longleftarrow$$P_{(k-gsize)}$$\times$((2$\times$(gsize-start)-count)/2$\times$(gc-start)+((cnt/2)$\times$(gc-start))$\times$ $P_{k}$
		\ENDFOR
   \ENDIF
\ENDFOR
\end{algorithmic}
\label{Algorithm:2}
\end{algorithm}

\begin{algorithm}[htp]
\caption{Random-fill}
\textbf{Input:} I(The image as 2D Matrix),gsize(Grid size).\par
\textbf{Output:} Modified image I after filling random color.\par
\textbf{Steps of the algorithm:}
\begin{algorithmic}[1]
\STATE nugb $\longleftarrow$ no of vertical ugb in image
\FOR{i $\longleftarrow$ 1 to nugb}
    \STATE p $\longleftarrow$ $i^{th}$ horizontal row of UGB's in I
    \FOR{ for each UGB in p}
    		\IF{the UGB is inside the isothetic cover}
    			\STATE fill the ugb with a random color
    		\ENDIF
	\ENDFOR
\ENDFOR
\IF{Condition}
    \STATE Do this
\ELSE
    \STATE Do that
\ENDIF
\end{algorithmic}
\label{Algorithm:2}
\end{algorithm}
\newpage






\section{EXPERIMENTAL RESULTS}
Results should be discussed here with appropriate comparison with other existing approaches (if applicable).




\newpage




\section{CONCLUSIONS}
State the concluding remarks and mention the future scopes and limitations of the project.




\newpage




\begin{thebibliography}{9}
	\bibitem{lamport94}
	Leslie Lamport,
	\emph{\LaTeX: a document preparation system},
	Addison Wesley, Massachusetts,
	2nd edition,
	1994.

	\bibitem{lamport95}
	Arindam Biswas, blabla.
	
\end{thebibliography}

\newpage


% BibTeX users use the following
% \bibliographystyle{plain}           % basic style
% \bibliography{Reference}            % name your BibTeX database file

% Non-BibTeX users use the following
%\begin{thebibliography}{}
%
% and use \bibitem to create references. Consult the Instructions
% for authors for reference list style.
%
%\bibitem{} Author. Title. Journal, Volume(Number):Pages, Year

\end{document}


